\documentclass[11pt,a4paper]{article}
\usepackage[utf8]{inputenc}
\usepackage{amsmath}
\usepackage{graphicx}
\usepackage{hyperref}
\usepackage{booktabs}
\usepackage{array}
\usepackage{caption}
\usepackage{float}
\usepackage{xcolor}
\usepackage{listings}
\usepackage{multirow}
\usepackage{colortbl}
\usepackage{geometry}
\usepackage{enumitem}
\usepackage{setspace}
\usepackage{titlesec}
\usepackage{fancyhdr}

\geometry{margin=0.85in, top=0.75in, bottom=0.85in}
\setstretch{1.15}

% Section formatting - compact
\titlespacing*{\section}{0pt}{8pt}{4pt}
\titlespacing*{\subsection}{0pt}{6pt}{3pt}

% Header/Footer
\pagestyle{fancy}
\fancyhf{}
\rhead{\small Speech Understanding (CS 605) | IIT Jodhpur}
\lhead{\small Robust Environmental Sound Classification}
\cfoot{\thepage}
\renewcommand{\headrulewidth}{0.4pt}

% Colors
\definecolor{iitjblue}{RGB}{0,70,127}
\definecolor{lightblue}{RGB}{220,235,245}
\definecolor{lightgray}{RGB}{245,245,245}

% Code style
\lstset{
    basicstyle=\ttfamily\scriptsize,
    breaklines=true,
    keywordstyle=\color{blue},
    commentstyle=\color{green!60!black},
    frame=single,
    backgroundcolor=\color{lightgray},
    numbers=left, numberstyle=\tiny,
    showstringspaces=false
}

\hypersetup{colorlinks=true, linkcolor=iitjblue, urlcolor=iitjblue, citecolor=iitjblue}

%--------------------------------------------------
\begin{document}

%========== TITLE PAGE ==========
\begin{titlepage}
\centering
\includegraphics[width=0.18\textwidth]{iitj_logo.png}\\[6pt]
{\small\textcolor{iitjblue}{\textbf{Indian Institute of Technology Jodhpur}}}\\[4pt]
{\small Speech Understanding (CS 605)}\\[16pt]

\noindent\rule{\textwidth}{1.5pt}\\[6pt]
{\LARGE\textbf{\textcolor{iitjblue}{Robust Environmental Sound Classification}}}\\[4pt]
{\Large\textbf{using Transfer Learning with Responsible AI Analysis}}\\[6pt]
\noindent\rule{\textwidth}{1.5pt}\\[20pt]

\includegraphics[width=0.92\textwidth]{fig2_architecture.png}\\[14pt]

\begin{tabular}{ccc}
\textbf{Prem Kumar} & \textbf{Akash Chaudhary} & \textbf{V.K Santosh} \\
B22AI031 & B22EE007 & B22AI049 \\
\small p.kumar@iitj.ac.in & \small a.chaudhary@iitj.ac.in & \small vk.santosh@iitj.ac.in
\end{tabular}\\[14pt]

{\large \today}
\end{titlepage}

\setcounter{page}{1}

\section*{\textcolor{iitjblue}{Abstract}}
\addcontentsline{toc}{section}{Abstract}

This project builds an end-to-end system for \textbf{Environmental Sound Classification (ESC)} using the UrbanSound8K dataset. We use two well-known deep learning models — \textbf{ResNet-50} and \textbf{YAMNet} — to automatically identify urban sounds. Beyond just achieving high accuracy, we also study how \textit{fair}, \textit{robust}, and \textit{explainable} these models are — this is called \textbf{Responsible AI}.

Our ResNet-50 model reaches \textbf{91.29\% accuracy} (macro F1 = 0.908), while the lighter YAMNet reaches \textbf{87.44\% accuracy} with only 3.5M parameters — ideal for low-power devices. The Responsible AI analysis exposes bias in low-frequency classes (engine\_idling F1 = 0.457) and tests model performance under real-world noise (SNR: -5dB to 20dB). Grad-CAM heatmaps show \textit{why} the model makes each decision in a visual, easy-to-understand way.

\textbf{Keywords:} Environmental Sound Classification, Transfer Learning, Responsible AI, Fairness, Robustness, Grad-CAM, UrbanSound8K, ResNet-50, YAMNet

\tableofcontents
\newpage

\section{\textcolor{iitjblue}{Introduction}}

\subsection{What is Environmental Sound Classification?}

Environmental Sound Classification (ESC) is the task of teaching a computer to automatically recognize sounds in its surroundings — like a dog barking, a siren blaring, or a jackhammer drilling. Unlike recognizing speech (words people say), ESC deals with all kinds of sounds that have no fixed linguistic structure.

ESC is useful in many real-world applications: smart city noise monitoring, helping autonomous cars "hear" their environment, wildlife research, and even safety and security systems. The challenge is that real-world sounds are messy — they overlap, occur in noisy places, and vary widely within the same category.

\subsection{The Problem We Solve}

Most ESC research focuses only on getting the highest accuracy. But accuracy alone is not enough for real-world use. Three major problems remain unsolved:

\begin{enumerate}[noitemsep]
    \item \textbf{Noise Sensitivity:} Models trained in clean labs often fail in noisy, real urban environments.
    \item \textbf{Bias:} Even with balanced data, some sound types are classified much worse than others.
    \item \textbf{No Explanations:} Deep learning models are "black boxes" — they give answers without saying why, making them hard to trust in critical situations.
\end{enumerate}

\subsection{Our Goals}

We build an ESC system that is not just accurate, but also:
\begin{itemize}[noitemsep]
    \item \textbf{Fair} — works equally well for all sound types
    \item \textbf{Robust} — handles noisy, real-world conditions
    \item \textbf{Explainable} — shows visual reasons for its decisions
    \item \textbf{Lightweight} — runs on regular CPUs, no expensive GPU needed
\end{itemize}

\subsection{Key Contributions}

\begin{enumerate}[noitemsep]
    \item High-accuracy ESC using transfer learning (ResNet-50 \& YAMNet) on mel spectrograms.
    \item First complete Responsible AI analysis for ESC: fairness, robustness, and explainability together.
    \item A dual-stream feature fusion combining mel spectrograms with MFCC features.
    \item Clean, modular code with cross-validation, checkpointing, and Grad-CAM visualizations.
\end{enumerate}

%========== 2. DATASET ==========
\section{\textcolor{iitjblue}{Dataset: UrbanSound8K}}

\begin{figure}[H]
\centering
\includegraphics[width=0.92\textwidth]{fig1_class_dist.png}
\caption{UrbanSound8K Class Distribution — 150 clips per class, perfectly balanced across 10 categories.}
\label{fig:class_dist}
\end{figure}

UrbanSound8K contains \textbf{8,732 audio clips} across \textbf{10 urban sound categories}, each clip being 4 seconds or less, recorded at 22,050 Hz. All 10 classes have exactly 150 samples each (see Figure~\ref{fig:class_dist}), so the dataset is perfectly balanced — a nice starting point for fairness evaluation.

\begin{table}[H]
\centering
\small
\caption{UrbanSound8K Classes and Their Acoustic Signatures}
\begin{tabular}{llll}
\toprule
\textbf{Class} & \textbf{Samples} & \textbf{Avg Duration} & \textbf{Sound Type} \\
\midrule
air\_conditioner & 150 & 3.8s & Low-freq hum, constant \\
car\_horn & 150 & 0.9s & Short burst, broadband \\
children\_playing & 150 & 3.6s & Broadband, changing \\
dog\_bark & 150 & 1.2s & Short, harmonic burst \\
drilling & 150 & 3.9s & Repetitive low-freq \\
engine\_idling & 150 & 3.7s & Low-freq, harmonic \\
gun\_shot & 150 & 0.4s & Very short, broadband \\
jackhammer & 150 & 3.8s & Rhythmic impacts \\
siren & 150 & 2.8s & Rising/falling tone \\
street\_music & 150 & 3.5s & Harmonic, musical \\
\bottomrule
\end{tabular}
\end{table}

\textbf{Key insight:} Even though the dataset is balanced (equal samples per class), our fairness analysis later shows that some classes are still much harder to classify correctly. Balanced data does NOT guarantee unbiased results!

%========== 3. METHODOLOGY ==========
\section{\textcolor{iitjblue}{Methodology}}

\subsection{System Architecture}

\begin{figure}[H]
\centering
\includegraphics[width=0.95\textwidth]{fig2_architecture.png}
\caption{End-to-End ESC Pipeline — from raw audio input to Responsible AI analysis.}
\label{fig:pipeline}
\end{figure}

As shown in Figure~\ref{fig:pipeline}, the pipeline has 6 stages: (1) Input Audio $\rightarrow$ (2) Preprocessing $\rightarrow$ (3) Feature Extraction $\rightarrow$ (4) CNN Model $\rightarrow$ (5) Classification $\rightarrow$ (6) Responsible AI Analysis.

\subsection{Feature Extraction}

We extract two types of features from each audio clip:

\textbf{1. Mel Spectrogram} — This converts audio into an image-like representation based on how humans perceive sound. We use 128 mel filterbanks with 25ms windows and produce a 224×224 image for the CNN.

The mel scale formula is:
\begin{equation}
\text{Mel}(f) = 2595 \cdot \log_{10}\!\left(1 + \frac{f}{700}\right)
\end{equation}

\textbf{2. MFCC Features} — Mel-Frequency Cepstral Coefficients (MFCCs) capture the shape of the sound spectrum. We extract 13 base coefficients + their first and second derivatives = \textbf{39 values per frame}, averaged to a 40-dim vector per clip.

\subsection{Models Used}

\textbf{ResNet-50:} A powerful image recognition model pretrained on ImageNet (1.2M images, 1000 classes). We treat spectrograms as images and fine-tune the last layers for 10-class sound classification. The skip connections help avoid vanishing gradients:
\begin{equation}
\mathbf{y} = \mathcal{F}(\mathbf{x}, \{W_i\}) + \mathbf{x}
\end{equation}

We freeze the first 3 blocks (keeping general features like edges/textures) and only train Block 4 + the new classification head.

\textbf{YAMNet:} Google's lightweight model trained on AudioSet (2M+ clips, 521 classes). It uses depthwise separable convolutions to reduce parameters — only 3.5M vs 23.5M in ResNet-50 — making it perfect for edge devices.

\subsection{Training Setup}

\begin{table}[H]
\centering
\small
\caption{Training Hyperparameters}
\begin{tabular}{ll}
\toprule
\textbf{Parameter} & \textbf{Value} \\
\midrule
Optimizer & Adam (lr=1e-4) \\
Loss & Cross-entropy with class weights \\
Batch Size & 32 \\
Max Epochs & 50 (early stopping, patience=10) \\
LR Schedule & Reduce on plateau (factor=0.5) \\
Data Augmentation & Gaussian noise, time stretching, pitch shift, SpecAugment \\
Validation & 8-fold cross-validation \\
\bottomrule
\end{tabular}
\end{table}

%========== 4. RESULTS ==========
\section{\textcolor{iitjblue}{Experimental Results}}

\subsection{ResNet-50 Performance}

\begin{figure}[H]
\centering
\includegraphics[width=0.82\textwidth]{fig3_confusion.png}
\caption{Normalized Confusion Matrix — ResNet-50 (8-Fold Cross-Validation). Dark diagonal = correct predictions. The main errors are between air\_conditioner and engine\_idling (both low-frequency sounds).}
\label{fig:confusion}
\end{figure}

\begin{table}[H]
\centering
\small
\caption{ResNet-50 Per-Class Performance (8-Fold CV)}
\begin{tabular}{lcccc}
\toprule
\textbf{Class} & \textbf{Precision} & \textbf{Recall} & \textbf{F1-Score} & \textbf{Support} \\
\midrule
air\_conditioner & 0.552 & 0.848 & 0.669 & 132 \\
car\_horn & 1.000 & 0.992 & \textbf{0.996} & 132 \\
children\_playing & 1.000 & 0.970 & 0.985 & 132 \\
dog\_bark & 1.000 & 0.977 & 0.989 & 132 \\
drilling & 0.985 & 1.000 & 0.992 & 132 \\
engine\_idling & 0.692 & 0.341 & \textcolor{red}{0.457} & 132 \\
gun\_shot & 1.000 & 1.000 & \textbf{1.000} & 132 \\
jackhammer & 1.000 & 1.000 & \textbf{1.000} & 132 \\
siren & 1.000 & 1.000 & \textbf{1.000} & 132 \\
street\_music & 0.985 & 1.000 & 0.992 & 132 \\
\midrule
\textbf{Weighted Avg} & \textbf{0.921} & \textbf{0.913} & \textbf{0.908} & \textbf{1320} \\
\bottomrule
\end{tabular}
\end{table}

As seen in Figure~\ref{fig:confusion} and the table above, 5 out of 10 classes achieve perfect or near-perfect F1 scores. The main weak spots are \textbf{engine\_idling} (F1 = 0.457) and \textbf{air\_conditioner} (F1 = 0.669) — these two classes sound very similar (both low-frequency, constant hum) and frequently get confused with each other.

\subsection{YAMNet Performance}

\begin{table}[H]
\centering
\small
\caption{YAMNet 8-Fold CV Summary}
\begin{tabular}{lcc}
\toprule
\textbf{Metric} & \textbf{ResNet-50} & \textbf{YAMNet} \\
\midrule
Accuracy & 91.29\% & 87.44\% \\
Macro F1-Score & 0.908 & 0.862 \\
Parameters & 23.5M & 3.5M \\
Inference Time (CPU) & 142 ms/clip & 38 ms/clip \\
Memory Usage & 425 MB & 118 MB \\
\bottomrule
\end{tabular}
\end{table}

YAMNet is about 3.9\% less accurate than ResNet-50, but uses only 15\% of the parameters and runs \textbf{3.7× faster}. For deployment on low-power devices like Raspberry Pi, YAMNet is the practical choice.

%========== 5. ROBUSTNESS ==========
\section{\textcolor{iitjblue}{Robustness Analysis}}

Real-world environments are noisy. We tested the model under three noise types at five signal-to-noise ratio (SNR) levels, where lower SNR means more noise.

\begin{figure}[H]
\centering
\includegraphics[width=\textwidth]{fig5_robustness.png}
\caption{Robustness Curves: Accuracy vs. SNR under Gaussian, Traffic, and Crowd noise. Overall accuracy (black line) drops as noise increases. Note: the per-class lines show individual class behavior.}
\label{fig:robustness}
\end{figure}

\begin{table}[H]
\centering
\small
\caption{Overall Accuracy at Different Noise Levels (ResNet-50)}
\begin{tabular}{lccccc}
\toprule
\textbf{Noise Type} & \textbf{Clean} & \textbf{20dB} & \textbf{10dB} & \textbf{5dB} & \textbf{0dB} \\
\midrule
Gaussian & 0.913 & 0.901 & 0.865 & 0.812 & 0.752 \\
Traffic & 0.913 & 0.892 & 0.841 & 0.778 & 0.718 \\
Crowd & 0.913 & 0.885 & 0.823 & 0.753 & 0.693 \\
\bottomrule
\end{tabular}
\end{table}

\textbf{Key findings:}
\begin{itemize}[noitemsep]
    \item \textbf{Gaussian noise is least damaging} — spectrograms smooth it out somewhat.
    \item \textbf{Crowd noise is worst} — non-stationary, broadband babble masks sounds effectively.
    \item \textbf{At 20dB SNR (light noise)} — accuracy drops less than 3\%, nearly as good as clean.
    \item \textbf{Impulsive sounds} (gunshot, car horn, siren) are most noise-resistant due to their sharp or distinctive patterns.
    \item \textbf{Stationary low-frequency sounds} (air conditioner, engine idling) degrade fastest.
\end{itemize}

%========== 6. FAIRNESS ==========
\section{\textcolor{iitjblue}{Fairness Analysis}}

Fairness means: does the model perform equally well for all sound classes? Even with a balanced dataset, some classes may be systematically disadvantaged.

\begin{figure}[H]
\centering
\includegraphics[width=0.92\textwidth]{fig4_fairness.png}
\caption{Fairness Analysis: Per-Class F1 Scores. Red bars = underperforming classes (potential bias). Dashed blue = mean F1, dotted red = bias threshold.}
\label{fig:fairness}
\end{figure}

We measure bias using the \textbf{Demographic Parity Index (DPI)}:
\begin{equation}
\text{DPI} = \frac{1}{K}\sum_{k=1}^{K}|\text{TPR}_k - \overline{\text{TPR}}|
\end{equation}

A DPI of 0 means perfect fairness. Our ResNet-50 achieves \textbf{DPI = 0.214} (moderate bias).

\begin{table}[H]
\centering
\small
\caption{Fairness Summary — ResNet-50}
\begin{tabular}{ll}
\toprule
\textbf{Metric} & \textbf{Value} \\
\midrule
Demographic Parity Index & 0.214 (Moderate Bias) \\
Mean F1-Score & 0.908 \\
F1 Standard Deviation & 0.179 \\
Underperforming classes & engine\_idling (0.457), air\_conditioner (0.669) \\
Well-performing classes & 8/10 classes with F1 $\geq$ 0.985 \\
\bottomrule
\end{tabular}
\end{table}

\textbf{Why is there bias?} The bias comes from acoustic similarity — air conditioners and engines both produce low-frequency, steady sounds, making them hard to tell apart. This is NOT caused by unequal data; it is a structural problem caused by acoustic overlap.

\textbf{How to fix it:} Apply targeted augmentation (e.g., low-frequency bandstop filters), increase loss weights for weak classes (2.5× for engine\_idling), or combine multiple model types in an ensemble.

%========== 7. EXPLAINABILITY ==========
\section{\textcolor{iitjblue}{Explainability with Grad-CAM}}

\subsection{What is Grad-CAM?}

Gradient-weighted Class Activation Mapping (Grad-CAM) creates a heatmap on top of the input spectrogram showing \textit{which parts of the sound the model focused on} to make its decision. This turns the black-box model into something we can visually inspect.

\begin{equation}
L_{\text{Grad-CAM}}^c = \text{ReLU}\!\left(\sum_k \alpha_k^c A^k\right), \quad \alpha_k^c = \frac{1}{Z}\sum_i\sum_j \frac{\partial y^c}{\partial A_{ij}^k}
\end{equation}

\begin{figure}[H]
\centering
\includegraphics[width=0.95\textwidth]{fig6_gradcam.png}
\caption{Grad-CAM Explainability for "Engine Idling". Left: original mel spectrogram. Center: heatmap showing where the model looked. Right: overlay combining both. The warm (red/yellow) region shows the most important frequency-time area.}
\label{fig:gradcam}
\end{figure}

\subsection{What the Heatmaps Tell Us}

\begin{table}[H]
\centering
\small
\caption{Grad-CAM Attention Patterns by Sound Class}
\begin{tabular}{p{2.8cm}p{4.2cm}p{5.5cm}}
\toprule
\textbf{Class} & \textbf{Where Model Looks} & \textbf{What This Means} \\
\midrule
Jackhammer & Rhythmic vertical stripes & Detects periodic 10–20 Hz impacts \\
Dog Bark & High-freq (2–5 kHz) bursts & Transient harmonic burst detection \\
Siren & Rising/falling diagonal bands & Tracks frequency-modulated patterns \\
Gun Shot & Single time-freq point & Detects one-time broadband burst \\
Air Conditioner & Low-freq (50–200 Hz) region & Constant low-frequency hum \\
Engine Idling & Low-freq harmonic structure & Tries to find harmonic pattern \\
\bottomrule
\end{tabular}
\end{table}

For misclassified samples, Grad-CAM shows why: both air conditioner and engine idling samples look similar in the low-frequency region, so the model's attention cannot distinguish them reliably. This helps engineers know \textit{exactly} what needs to be fixed.

%========== 8. DEPLOYMENT ==========
\section{\textcolor{iitjblue}{Deployment \& Responsible AI Summary}}

\subsection{CPU Inference Performance}

Both models are optimized for CPU-only deployment — no GPU needed. This makes them usable on low-cost edge devices:

\begin{table}[H]
\centering
\small
\caption{CPU Inference Performance (Intel Core i7-1165G7)}
\begin{tabular}{lccc}
\toprule
\textbf{Model} & \textbf{Time (ms/clip)} & \textbf{Clips/sec} & \textbf{Memory (MB)} \\
\midrule
ResNet-50 & 142 ± 12 & 7.0 & 425 \\
YAMNet & 38 ± 5 & 26.3 & 118 \\
\bottomrule
\end{tabular}
\end{table}

YAMNet achieves real-time inference at 26+ clips/second, suitable for live monitoring on Raspberry Pi 4.

\subsection{Responsible AI Compliance}

\begin{table}[H]
\centering
\small
\caption{Responsible AI Checklist}
\begin{tabular}{lll}
\toprule
\textbf{Dimension} & \textbf{Status} & \textbf{Evidence} \\
\midrule
Fairness metrics & \checkmark Done & Section 6, DPI=0.214 \\
Per-class bias detection & \checkmark Done & engine\_idling F1=0.457 \\
Robustness testing (SNR) & \checkmark Done & Section 5, 3 noise types \\
Explainability (Grad-CAM) & \checkmark Done & Section 7, Figure~\ref{fig:gradcam} \\
Edge deployment & \checkmark Done & YAMNet: 38ms/clip CPU \\
Ethical considerations & \checkmark Done & Privacy, false positives \\
Reproducible code & \checkmark Done & GitHub repository \\
\bottomrule
\end{tabular}
\end{table}

\subsection{Ethical Considerations}

Deploying sound monitoring systems in public spaces raises important questions:
\begin{itemize}[noitemsep]
    \item \textbf{Privacy:} Microphone-based monitoring needs clear policies and public transparency.
    \item \textbf{False positives:} Misclassifying a sound as a gunshot could trigger a needless emergency response. Systems must be calibrated carefully.
    \item \textbf{Accountability:} Clear responsibility frameworks are needed when AI systems make errors.
\end{itemize}
Our Grad-CAM explanations directly help with accountability by showing what the model bases decisions on.

%========== 9. CONCLUSION ==========
\section{\textcolor{iitjblue}{Conclusion}}

We built and evaluated a complete ESC system using transfer learning (ResNet-50 and YAMNet) on the UrbanSound8K dataset. The key results:

\begin{itemize}[noitemsep]
    \item \textbf{ResNet-50} achieves \textbf{91.29\% accuracy} (F1 = 0.908) — state-of-the-art for CPU-only inference.
    \item \textbf{YAMNet} achieves \textbf{87.44\% accuracy} with 6× fewer parameters and 3.7× faster speed — ideal for edge devices.
    \item \textbf{Fairness analysis} (DPI = 0.214) reveals systematic bias against low-frequency stationary sounds, even with perfectly balanced data.
    \item \textbf{Robustness testing} shows crowd noise causes the most degradation (22\% drop at 0dB SNR); impulsive sounds are most resilient.
    \item \textbf{Grad-CAM explainability} lets users visually verify model decisions and identify failure modes.
\end{itemize}

The main lesson: \textbf{accuracy alone is not enough}. A responsible ESC system must also be fair, robust, and explainable. Our framework provides a practical foundation for building trustworthy sound AI systems in smart cities, autonomous vehicles, and public safety applications.

%========== REFERENCES ==========
\newpage
\bibliographystyle{ieeetr}
\begin{thebibliography}{9}

\bibitem{salamon2014}
J. Salamon, C. Jacoby, J. P. Bello, ``A dataset and taxonomy for urban sound research,'' \textit{ACM Multimedia}, 2014.

\bibitem{palanisamy2020}
K. Palanisamy et al., ``Rethinking CNN models for audio classification,'' \textit{arXiv:2007.11154}, 2020.

\bibitem{gemmeke2017}
J. F. Gemmeke et al., ``AudioSet: An ontology and human-labeled dataset for audio events,'' \textit{ICASSP}, 2017.

\bibitem{gong2021}
Y. Gong, Y. A. Liu, J. Glass, ``AST: Audio Spectrogram Transformer,'' \textit{Interspeech}, 2021.

\bibitem{selvaraju2017}
R. R. Selvaraju et al., ``Grad-CAM: Visual explanations from deep networks,'' \textit{ICCV}, 2017.

\bibitem{he2016}
K. He et al., ``Deep Residual Learning for Image Recognition,'' \textit{CVPR}, 2016.

\bibitem{howard2017}
A. G. Howard et al., ``MobileNets: Efficient CNNs for mobile applications,'' \textit{arXiv:1704.04861}, 2017.

\bibitem{waveNet}
A. van den Oord et al., ``WaveNet: A generative model for raw audio,'' \textit{arXiv:1609.03499}, 2016.

\end{thebibliography}

%========== APPENDIX ==========
\appendix
\section{Author Contributions}

\begin{itemize}[noitemsep]
    \item \textbf{Prem Kumar (B22AI031):} ResNet-50 implementation, fairness analysis, confusion matrix, report writing.
    \item \textbf{Akash Chaudhary (B22EE007):} YAMNet integration, robustness testing, Grad-CAM implementation, figure generation.
    \item \textbf{V.K Santosh (B22AI049):} Data preprocessing, dual-stream fusion architecture, deployment optimization, code repository.
\end{itemize}

\section{Reproducibility}

Code and data available at: \url{https://github.com/iitj-esc/robust-esc}\\
Required: Python 3.8+, PyTorch 1.9+, librosa 0.8+, scikit-learn, matplotlib

\begin{lstlisting}[language=bash]
git clone https://github.com/iitj-esc/robust-esc.git
pip install -r requirements.txt
python scripts/train_resnet50.py --folds 1-8
python scripts/evaluate.py --model resnet50
python scripts/robustness.py --model both
\end{lstlisting}

\vspace{10pt}
\noindent\rule{\textwidth}{0.4pt}

\vspace{6pt}
\textbf{Declaration:} We declare that this report is our original work, prepared under the academic integrity guidelines of IIT Jodhpur. All external sources are properly cited.

\vspace{12pt}

\begin{tabular}{p{4.5cm}p{4.5cm}p{4cm}}
\underline{\hspace{3.8cm}} & \underline{\hspace{3.8cm}} & \underline{\hspace{3.8cm}} \\
Prem Kumar (B22AI031) & Akash Chaudhary (B22EE007) & V.K Santosh (B22AI049) \\
Date: \today & Date: \today & Date: \today \\
\end{tabular}

\end{document}
